\section{Step 10: Final Robustness Check with Simpler Fixed-Effects and Cluster-Robust Regression}
\label{sec:step10_final_robustness}

\subsection{Objective}
The objective of Step~10 is to perform the final robustness check of the frontal condition effect
using a simpler regression framework than the mixed-effects models used in Steps~6--9.

This step is intended to answer one last question:
\begin{quote}
If the same frontal outcome is analyzed with a simpler fixed-effects regression and robust standard errors,
does the abstract-versus-concrete effect remain, weaken further, or disappear?
\end{quote}

This step should be the \textbf{last reanalysis step} before write-up.
Its purpose is not to open a new inferential branch, but to determine whether the remaining frontal
tendency is robust to a simpler and more transparent specification.

\subsection{Why Step~10 Is Needed}
The previous steps showed a consistent but fragile pattern:
\begin{itemize}
    \item the Step~5 fixed-window front ROI result was near threshold;
    \item the Step~6 covariate-adjusted front ROI model was even closer to threshold;
    \item the Step~7 duration-aware front ROI model remained positive but non-significant;
    \item the Step~9 block-adjusted front ROI model weakened further and became clearly non-significant.
\end{itemize}

At the same time, the mixed-effects reports repeatedly indicated that:
\begin{itemize}
    \item the optimizer converged with boundary warnings;
    \item participant random-intercept variance was extremely small;
    \item question random-intercept variance was extremely small.
\end{itemize}

This suggests that the data may be close to a simpler regression limit.
Therefore, Step~10 should test whether the main conclusion changes when the mixed-effects machinery
is replaced by a simpler fixed-effects and robust-inference approach.

\subsection{Core Analytical Decision}
To keep Step~10 interpretable, the following elements should remain fixed from Step~9:
\begin{enumerate}
    \item the same Step~9 frontal outcome;
    \item the same trial table and block metadata;
    \item the same included participant-session cohort in the primary analysis;
    \item the same item covariates:
    \begin{itemize}
        \item standardized log total word count;
        \item standardized ENEM item correctness.
    \end{itemize}
\end{enumerate}

The main change in Step~10 is:
\begin{quote}
replace the mixed-effects model with a simpler regression using participant fixed effects and
cluster-robust standard errors.
\end{quote}

\subsection{Primary Outcome}
The primary Step~10 outcome should remain the same as in Step~9:
\begin{quote}
the trial-level left frontal ROI HbO response computed from the duration-aware variable window.
\end{quote}

Keeping the same dependent variable ensures that Step~10 isolates the effect of changing the
modeling framework rather than changing the signal summary.

\subsection{Primary ROI Definition}
The primary ROI remains the same left frontal ROI used in Steps~5--9.

\subsubsection{Front ROI HbO channels}
\begin{verbatim}
S2_D6 hbo
S2_D4 hbo
S1_D6 hbo
S1_D3 hbo
S5_D6 hbo
S5_D4 hbo
S5_D3 hbo
\end{verbatim}

\subsubsection{Short-separation pairs excluded from inference}
\begin{verbatim}
S1_D8
S2_D9
S3_D10
S4_D11
S5_D12
S6_D13
S7_D14
S8_D15
\end{verbatim}

\subsection{Important Identifiability Note}
A full model with both \textbf{participant fixed effects} and \textbf{question fixed effects}
is \textbf{not suitable} for the primary Step~10 condition test.

The reason is that:
\begin{itemize}
    \item condition is a question-level attribute;
    \item word count is a question-level attribute;
    \item ENEM item correctness is a question-level attribute.
\end{itemize}

If question fixed effects are added, these variables become perfectly collinear with the question dummies
and their coefficients cannot be estimated.

Therefore, the primary Step~10 model should use:
\begin{itemize}
    \item participant fixed effects;
    \item question-level covariates;
    \item cluster-robust standard errors.
\end{itemize}

Question-level dependence is then handled through robust inference and influence diagnostics rather than
through question fixed effects.

\subsection{Variables Used in the Primary Model}
Let:
\begin{equation}
y^{(F)}_{ik}
\end{equation}
denote the Step~9 duration-aware front ROI HbO response for participant \(i\) on trial \(k\).

Let:
\begin{equation}
A_{ik} =
\begin{cases}
1 & \text{if the trial is Abstract} \\
0 & \text{if the trial is Concrete}
\end{cases}
\end{equation}

Let:
\begin{equation}
g_{ik} \in \{1,\dots,10\}
\end{equation}
be realized global block order, and
\begin{equation}
p_{ik} \in \{1,2,3\}
\end{equation}
be within-block position.

Let:
\begin{equation}
w_k = z\!\left(\log(\texttt{total\_word\_count}_k)\right)
\end{equation}
and
\begin{equation}
c_k = z(\texttt{correctness}_k).
\end{equation}

\subsection{Primary Step~10 Model}
The primary Step~10 model should be:
\begin{equation}
y^{(F)}_{ik}
=
\beta_0
+
\beta_1 A_{ik}
+
\beta_2 g_{ik}
+
\beta_3 p_{ik}
+
\beta_4 w_k
+
\beta_5 c_k
+
\alpha_i
+
\varepsilon_{ik},
\end{equation}
where:
\begin{itemize}
    \item \(\alpha_i\) is a participant fixed effect;
    \item \(\varepsilon_{ik}\) is the residual term.
\end{itemize}

This model keeps participant-level heterogeneity explicit without relying on a random-effects variance
that appears to be close to zero in the earlier mixed models.

\subsection{Primary Inference Strategy}
The primary inferential target remains:
\begin{equation}
\beta_1,
\end{equation}
the adjusted abstract-versus-concrete effect in the left frontal ROI.

The preferred inference method should be:
\begin{quote}
two-way cluster-robust standard errors by participant and question
\end{quote}
if the software implementation is stable.

If two-way clustered standard errors are not available or are unstable, the primary fallback should be:
\begin{quote}
participant-clustered standard errors,
\end{quote}
with question-clustered standard errors reported as a sensitivity analysis.

\subsection{Primary Hypothesis}
The primary hypotheses are:

\paragraph{Null hypothesis}
\begin{equation}
H_0: \beta_1 = 0
\end{equation}

\paragraph{Alternative hypothesis}
\begin{equation}
H_1: \beta_1 \neq 0
\end{equation}

The significance threshold should remain:
\begin{verbatim}
alpha = 0.05
\end{verbatim}

\subsection{Why Condition Order Is Not in the Primary Model}
Participant-level condition order was included in Step~9 as a mixed-model covariate.
In Step~10, participant fixed effects absorb all participant-level constants.
Therefore, condition order cannot be estimated separately in the primary fixed-effects regression.

This is acceptable, because the point of Step~10 is not to rebuild every Step~9 coefficient.
The point is to test whether the frontal condition effect survives in a simpler and more robust
participant-fixed-effects framework.

\subsection{Secondary Sensitivity Models}
To complete the final robustness check, Step~10 should include the following secondary analyses.

\subsubsection{Sensitivity Model A: Participant-clustered standard errors}
Fit the primary model and report participant-clustered standard errors.

\subsubsection{Sensitivity Model B: Question-clustered standard errors}
Fit the primary model and report question-clustered standard errors.

\subsubsection{Sensitivity Model C: HC3 heteroskedasticity-robust standard errors}
Fit the primary model and report HC3 robust standard errors.

\subsubsection{Sensitivity Model D: Step~6 fixed-window benchmark}
Repeat the same regression structure using the Step~6 fixed-window front ROI outcome instead of the
Step~9 duration-aware outcome.

This provides a final benchmark between:
\begin{itemize}
    \item the duration-aware frontal outcome;
    \item the fixed-window frontal outcome.
\end{itemize}

\subsubsection{Sensitivity Model E: No-override cohort}
Repeat the primary regression after removing the two inherited manual-inclusion sessions:
\begin{itemize}
    \item PID025 / 2025-11-25 002
    \item PID034 / 2025-11-27 003
\end{itemize}

This is important because those sessions were retained by explicit analyst override in the earlier
pipeline and remained part of the downstream cohort.

\subsection{Leave-One-Question-Out Influence Analysis}
Because question fixed effects cannot be used in the primary condition model, Step~10 should include
a leave-one-question-out influence analysis.

For each question \(q\):
\begin{itemize}
    \item remove all trials corresponding to \(q\);
    \item refit the primary participant-fixed-effects model;
    \item store the new condition coefficient \(\hat{\beta}_1^{(-q)}\);
    \item store the new p-value and confidence interval.
\end{itemize}

This analysis will show whether the frontal condition effect depends heavily on one or a few questions.

\subsection{Optional Leave-One-Participant-Out Check}
If computationally feasible, Step~10 may also include a leave-one-participant-out analysis:
\begin{itemize}
    \item remove one participant-session at a time;
    \item refit the primary model;
    \item record the condition coefficient and p-value.
\end{itemize}

This is optional, but useful for verifying that the result is not being driven by one unusually
influential participant.

\subsection{Secondary Back ROI Benchmark}
The back ROI has been consistently null and should remain secondary.
Still, one benchmark participant-fixed-effects model may be fit for the back ROI:
\begin{equation}
y^{(B)}_{ik}
=
\beta_0
+
\beta_1 A_{ik}
+
\beta_2 g_{ik}
+
\beta_3 p_{ik}
+
\beta_4 w_k
+
\beta_5 c_k
+
\alpha_i
+
\varepsilon_{ik}.
\end{equation}

This model is included only as a final benchmark and should not be emphasized.

\subsection{Step-by-Step Workflow}

\subsubsection{Step~10.1: Freeze the Step~9 trial table}
Start from the same Step~9 trial-level duration-aware frontal outcome and the same merged block metadata.

\subsubsection{Step~10.2: Build the primary regression table}
Create one row per valid trial, including:
\begin{itemize}
    \item participant ID;
    \item question ID;
    \item condition;
    \item duration-aware front ROI outcome;
    \item global block order;
    \item within-block position;
    \item standardized log total word count;
    \item standardized ENEM item correctness.
\end{itemize}

\subsubsection{Step~10.3: Fit the participant-fixed-effects primary model}
Fit the primary OLS regression with participant fixed effects and robust inference.

\subsubsection{Step~10.4: Compute the primary robust standard errors}
Use:
\begin{itemize}
    \item two-way clustered standard errors by participant and question, if available;
    \item otherwise participant-clustered standard errors as the primary fallback.
\end{itemize}

\subsubsection{Step~10.5: Fit the alternative robust-SE versions}
Refit the same model with:
\begin{itemize}
    \item participant-clustered SE;
    \item question-clustered SE;
    \item HC3 SE.
\end{itemize}

\subsubsection{Step~10.6: Fit the fixed-window benchmark}
Repeat the same participant-fixed-effects regression using the Step~6 fixed-window frontal outcome.

\subsubsection{Step~10.7: Fit the no-override sensitivity}
Repeat the primary duration-aware regression after removing PID025 and PID034.

\subsubsection{Step~10.8: Run the leave-one-question-out analysis}
Refit the primary model 30 times, each time dropping one question.

\subsubsection{Step~10.9: Optionally run the leave-one-participant-out analysis}
Refit the primary model once per participant-session exclusion, if desired.

\subsubsection{Step~10.10: Fit the secondary back ROI benchmark}
Fit the back ROI participant-fixed-effects benchmark model.

\subsection{Primary Reporting Quantities}
The primary Step~10 results table must report:
\begin{itemize}
    \item number of included participants;
    \item number of included questions;
    \item number of included trials;
    \item condition coefficient \(\beta_1\);
    \item standard error;
    \item test statistic;
    \item p-value;
    \item 95\% confidence interval;
    \item block-order coefficient;
    \item within-block-position coefficient;
    \item log-word-count coefficient;
    \item ENEM-correctness coefficient;
    \item robust-SE type used in the primary report.
\end{itemize}

\subsection{Secondary Reporting Quantities}
The secondary Step~10 tables should report:
\begin{itemize}
    \item participant-clustered results;
    \item question-clustered results;
    \item HC3 results;
    \item fixed-window benchmark results;
    \item no-override sensitivity results;
    \item leave-one-question-out coefficient range;
    \item optional leave-one-participant-out coefficient range;
    \item back-ROI benchmark results.
\end{itemize}

\subsection{Primary Conclusion Rule}
The final Step~10 conclusion should be based on the condition coefficient \(\beta_1\) from the
primary participant-fixed-effects model.

The interpretation rule should be:
\begin{itemize}
    \item if the coefficient remains significant under the primary robust inference, conclude that
    the frontal condition effect survives a simpler fixed-effects robustness check;
    \item if the coefficient is non-significant and remains unstable across the sensitivity models,
    conclude that the residual frontal tendency does not survive a final simpler-model robustness analysis;
    \item if the coefficient changes sign or depends strongly on a small number of questions or participants,
    conclude that the frontal tendency is fragile and highly influence-sensitive.
\end{itemize}

\subsection{Interpretation Templates}

\paragraph{Template if the frontal condition effect survives.}
In the final participant-fixed-effects robustness analysis, the abstract-versus-concrete effect
in the left frontal ROI remained after controlling for block order, within-block position,
question length, and item difficulty
(\(\beta=\ldots\), \(SE=\ldots\), \(p=\ldots\), 95\% CI \([\ldots,\ldots]\)).
This suggests that the frontal condition tendency is not an artifact of the mixed-effects specification.

\paragraph{Template if the frontal condition effect does not survive.}
In the final participant-fixed-effects robustness analysis, the abstract-versus-concrete effect
in the left frontal ROI was not statistically significant
(\(\beta=\ldots\), \(SE=\ldots\), \(p=\ldots\), 95\% CI \([\ldots,\ldots]\)).
Together with the previous block-adjusted and duration-aware results, this suggests that the frontal
condition tendency remains descriptive but does not receive robust statistical support in the current dataset.

\paragraph{Template if the result is highly influence-sensitive.}
In the final participant-fixed-effects robustness analysis, the abstract-versus-concrete coefficient
was unstable across influence checks, indicating that the apparent frontal tendency depends strongly
on a small subset of questions and/or participants.
This pattern argues against a robust condition-level frontal effect in the current dataset.

\subsection{Expected Outputs}

\subsubsection*{Tables}
\begin{verbatim}
01_step10_primary_fe_table.csv
02_step10_primary_fe_results.csv
03_step10_cluster_participant_results.csv
04_step10_cluster_question_results.csv
05_step10_hc3_results.csv
06_step10_fixed_window_benchmark.csv
07_step10_no_override_results.csv
08_step10_leave_one_question_out.csv
09_step10_leave_one_participant_out.csv
10_step10_back_roi_benchmark.csv
11_step10_exclusion_log.csv
12_step10_final_conclusion.csv
\end{verbatim}

\subsubsection*{Figures}
\begin{verbatim}
figures/step10/primary_condition_coefficient_plot.png
figures/step10/robust_se_comparison.png
figures/step10/fixed_vs_variable_fe_benchmark.png
figures/step10/no_override_comparison.png
figures/step10/leave_one_question_out_plot.png
figures/step10/leave_one_participant_out_plot.png
\end{verbatim}

\subsubsection*{Reports}
\begin{verbatim}
reports/step10/step10_final_robustness_report.tex
reports/step10/tables/primary_fe_table.tex
reports/step10/tables/sensitivity_table.tex
reports/step10/text/final_conclusion.tex
\end{verbatim}

\subsection{Minimal Figures for the Main Report}
The main Step~10 report should include at least:
\begin{enumerate}
    \item coefficient plot for the primary participant-fixed-effects model;
    \item comparison of robust-SE variants;
    \item fixed-window versus duration-aware benchmark comparison;
    \item no-override sensitivity comparison;
    \item leave-one-question-out condition coefficient plot.
\end{enumerate}

\subsection{Quality-Control Checks}
Before Step~10 is considered complete, verify that:
\begin{enumerate}
    \item the Step~9 trial-level frontal outcome is unchanged;
    \item participant fixed effects are correctly constructed;
    \item question-level covariates are correctly merged;
    \item question fixed effects are not mistakenly included in the primary condition model;
    \item the primary robust-SE specification is documented clearly;
    \item the no-override sensitivity removes only PID025 and PID034;
    \item leave-one-question-out runs cover all included questions;
    \item all sensitivity models are clearly labeled as secondary.
\end{enumerate}

\subsection{Minimal Completion Criteria}
Step~10 is complete when:
\begin{itemize}
    \item the primary participant-fixed-effects model has been fit;
    \item robust-SE sensitivity analyses have been completed;
    \item the fixed-window benchmark has been fit;
    \item the no-override sensitivity has been fit;
    \item the leave-one-question-out analysis has been completed;
    \item the back-ROI benchmark has been completed;
    \item the final robustness conclusion has been written.
\end{itemize}

\subsection{Short Practical Summary}
In simple terms, Step~10 should:
\begin{quote}
take the Step~9 duration-aware frontal outcome, replace the mixed model with a simpler
participant-fixed-effects regression, use robust standard errors, test whether the frontal
condition effect survives, repeat the analysis without the two inherited manual-override sessions,
check whether the result depends heavily on one or a few questions, and then stop the reanalysis cycle.
\end{quote}